\section{METHODOLOGY}

\subsection{Methodology Overview}
\begin{frame}[t]{Methodology Overview}
    \begin{block}{Methodology Direction}
        The study follows a design and development methodology: review theory, design reusable architecture, implement the framework, and validate it through experiments.
    \end{block}
    % overall architecture workflow
    \begin{figure}[H]
        \centering
        \resizebox{\textwidth}{!}{%
            \begin{tikzpicture}[
                    node distance=2.2cm,
                    every node/.style={
                            draw,
                            fill=lightblue,
                            align=center,
                            rounded corners,
                            thick,
                            font=\large,
                            minimum width=4.2cm,
                            minimum height=2.0cm
                        }
                ]
                \node (r1) {Literature Review\\and Requirements Analysis};
                \node (r2) [right=of r1] {Software\\Architecture Design};
                \node (r3) [right=of r2] {Framework\\Construction};
                \node (r4) [right=of r3] {Experimental\\Validation};
                \node (r5) [right=of r4] {Deployment\\PyPi Package};
                \node (r6) [below=of r5, fill=mainblue, text=white] {Documentation};

                \draw[->] (r1) -- (r2);
                \draw[->] (r2) -- (r3);
                \draw[->] (r3) -- (r4);
                \draw[->] (r4) -- (r5);
                \draw[->] (r5) -- (r6);
            \end{tikzpicture}%
        }
    \end{figure}
\end{frame}

\subsection{Implementation Design Concept}
\begin{frame}[fragile,t]{Implementation Design Concept}
    \centering
    \scriptsize
    \vspace{-0.25cm}

    \resizebox{0.95\textwidth}{!}{%
        \begin{tikzpicture}

            \node[parambox] (param) at (0,0) {
                User Parameter\\
                \texttt{"mean"}, \texttt{"rbf"}, \texttt{"mse"}
            };

            \node[factorybox] (factory) at (3.6,0) {
                Factory\\
                Selects component
            };

            \node[modulebox] (module) at (7.6,0) {
                Concrete Module\\
                Actual implementation
            };

            \node[basebox] (base) at (12.4,0) {
                Base Class\\
                Common interface
            };

            \draw[archarrow] (param) -- (factory);
            \draw[archarrow] (factory) -- (module);
            \draw[archdashed] (module) --
            node[midway, above, font=\scriptsize, fill=white, inner sep=1pt] {inherits from}
            (base);

            \node[factorybox] (factoryex) at (3.6,-1.6) {
                \texttt{CombinerFactory}\\
                \texttt{KernelFactory}\\
                \texttt{LossFactory}
            };

            \node[modulebox] (moduleex) at (7.7,-1.6) {
                \texttt{MeanCombiner}\\
                \texttt{RBFKernel}\\
                \texttt{MSELoss}
            };

            \node[basebox] (baseex) at (12.4,-1.6) {
                \texttt{BaseCombiner}\\
                \texttt{BaseKernel}\\
                \texttt{BaseLoss}
            };

            \draw[archarrow] (factory) -- (factoryex);
            \draw[archarrow] (factoryex) -- (moduleex);
            \draw[archdashed] (moduleex) --
            node[midway, above, font=\scriptsize, fill=white, inner sep=1pt] {inherits from}
            (baseex);

        \end{tikzpicture}
    }

    \vspace{0.25cm}

    \begin{block}{Purpose}
        The user provides simple parameter names, factories create the correct concrete modules, and base classes keep all components consistent and extensible.
    \end{block}

\end{frame}

\subsection{KFCProcedure Architecture and Usage}

\begin{frame}[fragile,t]{KFCProcedure Architecture}
    \centering
    \scriptsize
    \vspace{-0.25cm}

    \resizebox{0.96\textwidth}{!}{%
        \begin{tikzpicture}

            \node[mainbox] (kfc) at (0,0) {
                \texttt{KFCProcedure}
            };

            % Row 1
            \node[stepbox] (kstep) at (3.1,1.35) {
                K-Step
            };

            \node[factorybox] (kfactory) at (6.3,1.35) {
                Divergence\\
                Factory
            };

            \node[modulebox] (kmodule) at (9.7,1.35) {
                Divergence\\
                Module
            };

            \node[basebox] (kbase) at (13.1,1.35) {
                Base\\
                Divergence
            };

            % Row 2
            \node[stepbox] (fstep) at (3.1,0) {
                F-Step
            };

            \node[factorybox] (ffactory) at (6.3,0) {
                Model\\
                Factory
            };

            \node[modulebox] (fmodule) at (9.7,0) {
                Local Model\\
                Module
            };

            \node[basebox] (fbase) at (13.1,0) {
                Base\\
                Local Model
            };

            % Row 3
            \node[stepbox] (cstep) at (3.1,-1.35) {
                C-Step
            };

            \node[factorybox] (cfactory) at (6.3,-1.35) {
                Combiner\\
                Factory
            };

            \node[modulebox] (cmodule) at (9.7,-1.35) {
                Combiner\\
                Module
            };

            \node[basebox] (cbase) at (13.1,-1.35) {
                Base\\
                Combiner
            };

            % Arrows from main
            \draw[archarrow] (kfc) -- (kstep);
            \draw[archarrow] (kfc) -- (fstep);
            \draw[archarrow] (kfc) -- (cstep);

            % Row arrows
            \draw[archarrow] (kstep) -- (kfactory);
            \draw[archarrow] (kfactory) -- (kmodule);
            \draw[archdashed] (kmodule) -- (kbase);

            \draw[archarrow] (fstep) -- (ffactory);
            \draw[archarrow] (ffactory) -- (fmodule);
            \draw[archdashed] (fmodule) -- (fbase);

            \draw[archarrow] (cstep) -- (cfactory);
            \draw[archarrow] (cfactory) -- (cmodule);
            \draw[archdashed] (cmodule) -- (cbase);

        \end{tikzpicture}
    }

    \vspace{0.20cm}

    \begin{block}{Implementation Idea}
        Each KFC stage uses a factory to select a concrete module, while the base class defines a consistent interface for extension.
    \end{block}

\end{frame}

\begin{frame}[fragile, t]{KFCProcedure Usage}
    \vspace{-0.5cm}
    \begin{table}
        \scriptsize
        \renewcommand{\arraystretch}{1.2}
        \setlength{\tabcolsep}{1.8pt}
        \begin{tabularx}{\textwidth}{
                >{\raggedright\arraybackslash}m{2.55cm}
                >{\raggedright\arraybackslash}m{2.55cm}
                >{\raggedright\arraybackslash}X}
            \toprule
            \textbf{Parameter}          &
            \textbf{Required / Default} &
            \textbf{Purpose}                                                                 \\
            \midrule
            \texttt{divergences}        &
            \texttt{yes}                &
            List of Bregman divergences strings or instances for clustering.                 \\
            \texttt{local\_model}       &
            \texttt{yes}                &
            Base learner string or instance used for cluster-wise local modeling.            \\
            \texttt{combiner}           &
            \texttt{yes}                &
            Strategy algorithm handling ensemble aggregation.                                \\
            \texttt{task}               &
            \texttt{"regression"}       &
            Type of learning objective (\texttt{"regression"} or \texttt{"classification"}). \\
            \texttt{n\_clusters}        &
            \texttt{3}                  &
            Number of targeted clusters per divergence metric.                               \\
            \texttt{max\_iter}          &
            \texttt{300}                &
            Maximum iterations allowed for clustering optimization.                          \\
            \texttt{tol}                &
            \texttt{1e-4}               &
            Convergence numerical tolerance threshold.                                       \\
            \texttt{verbose}            &
            \texttt{0}                  &
            Debug logging depth configuration level (0 to 3).                                \\
            \texttt{random\_state}      &
            \texttt{None}               &
            Reproducibility random seed anchor.                                              \\
            \bottomrule
        \end{tabularx}
    \end{table}
    \tiny *Note: Additional parameters can be passed explicitly via \texttt{divergences\_params}, \texttt{local\_model\_params}, and \texttt{combiner\_params} dictionaries.
\end{frame}

\begin{frame}[fragile, t]{KFCProcedure Usage}
    \begin{minted}[
        fontsize=\scriptsize,
        breaklines,
        linenos,
        frame=single,
        framesep=2mm
    ]{python}
        from kfc_procedure import KFCProcedure

        model = KFCProcedure(
            divergences=["euclidean"], # euclidean, gkl, is, logistic
            local_model="lasso",       # linear_regression, ridge, lasso
            combiner="mean",           # mean
            task="regression",         # regression or classification
            n_clusters=3,              # number of clusters per divergence
            max_iter=300,              # maximum iterations for clustering
            tol=1e-4,                  # convergence numerical tolerance threshold.
            verbose=1,                 # debug logging level (0 to 3)
            random_state=42            # reproducibility random seed
        )
        model.fit(X_train, y_train)
        preds = model.predict(X_test)

    \end{minted}
\end{frame}

\subsection{COBRA-Based Aggregation}

\begin{frame}[fragile,t]{GradientCOBRA Implementation Architecture}
    \centering
    \scriptsize
    \vspace{-0.25cm}

    \resizebox{0.96\textwidth}{!}{%
        \begin{tikzpicture}

            \node[mainbox] (gcobra) at (0,0) {
                \texttt{GradientCOBRA}
            };

            % Parameter layer
            \node[parambox] (estimators) at (3.2,2.2) {\texttt{estimators}};
            \node[parambox] (distance)   at (3.2,1.35) {\texttt{distance}};
            \node[parambox] (kernel)     at (3.2,0.50) {\texttt{kernel}};
            \node[parambox] (optimizer)  at (3.2,-0.35) {\texttt{optimizer}};
            \node[parambox] (loss)       at (3.2,-1.20) {\texttt{loss}};
            \node[parambox] (aggregator) at (3.2,-2.05) {\texttt{aggregator}};

            % Factory layer
            \node[factorybox] (estfactory)    at (6.5,2.2) {Estimator\\Factory};
            \node[factorybox] (distfactory)   at (6.5,1.35) {Distance\\Factory};
            \node[factorybox] (kernelfactory) at (6.5,0.50) {Kernel\\Factory};
            \node[factorybox] (optfactory)    at (6.5,-0.35) {Optimizer\\Factory};
            \node[factorybox] (lossfactory)   at (6.5,-1.20) {Loss\\Factory};
            \node[factorybox] (aggfactory)    at (6.5,-2.05) {Aggregator\\Factory};

            % Concrete module layer
            \node[modulebox] (estmodule)    at (9.8,2.2) {Estimator\\Module};
            \node[modulebox] (distmodule)   at (9.8,1.35) {Distance\\Module};
            \node[modulebox] (kernelmodule) at (9.8,0.50) {Kernel\\Module};
            \node[modulebox] (optmodule)    at (9.8,-0.35) {Optimizer\\Module};
            \node[modulebox] (lossmodule)   at (9.8,-1.20) {Loss\\Module};
            \node[modulebox] (aggmodule)    at (9.8,-2.05) {Aggregator\\Module};

            % Base layer
            \node[basebox] (baseest)    at (13.1,2.2) {Base\\Estimator};
            \node[basebox] (basedist)   at (13.1,1.35) {Base\\Distance};
            \node[basebox] (basekernel) at (13.1,0.50) {Base\\Kernel};
            \node[basebox] (baseopt)    at (13.1,-0.35) {Base\\Optimizer};
            \node[basebox] (baseloss)   at (13.1,-1.20) {Base\\Loss};
            \node[basebox] (baseagg)    at (13.1,-2.05) {Base\\Aggregator};

            \node[outputbox] (output) at (16.4,0) {
                Final Prediction\\
                \texttt{predict()}
            };

            % Bus coordinates
            \coordinate (leftbus) at (1.75,0);
            \coordinate (rightbus) at (14.65,0);

            % Left bus
            \draw[archarrow] (gcobra) -- (leftbus);
            \draw[thick] (leftbus |- estimators) -- (leftbus |- aggregator);

            \draw[archarrow] (leftbus |- estimators) -- (estimators);
            \draw[archarrow] (leftbus |- distance) -- (distance);
            \draw[archarrow] (leftbus |- kernel) -- (kernel);
            \draw[archarrow] (leftbus |- optimizer) -- (optimizer);
            \draw[archarrow] (leftbus |- loss) -- (loss);
            \draw[archarrow] (leftbus |- aggregator) -- (aggregator);

            % Parameter to factory
            \draw[archarrow] (estimators) -- (estfactory);
            \draw[archarrow] (distance) -- (distfactory);
            \draw[archarrow] (kernel) -- (kernelfactory);
            \draw[archarrow] (optimizer) -- (optfactory);
            \draw[archarrow] (loss) -- (lossfactory);
            \draw[archarrow] (aggregator) -- (aggfactory);

            % Factory to module
            \draw[archarrow] (estfactory) -- (estmodule);
            \draw[archarrow] (distfactory) -- (distmodule);
            \draw[archarrow] (kernelfactory) -- (kernelmodule);
            \draw[archarrow] (optfactory) -- (optmodule);
            \draw[archarrow] (lossfactory) -- (lossmodule);
            \draw[archarrow] (aggfactory) -- (aggmodule);

            % Module inherits from base
            \draw[archdashed] (estmodule) -- (baseest);
            \draw[archdashed] (distmodule) -- (basedist);
            \draw[archdashed] (kernelmodule) -- (basekernel);
            \draw[archdashed] (optmodule) -- (baseopt);
            \draw[archdashed] (lossmodule) -- (baseloss);
            \draw[archdashed] (aggmodule) -- (baseagg);

            % Right bus to output
            \draw[thick] (rightbus |- baseest) -- (rightbus |- baseagg);

            \draw[archarrow] (baseest) -- (rightbus |- baseest);
            \draw[archarrow] (basedist) -- (rightbus |- basedist);
            \draw[archarrow] (basekernel) -- (rightbus |- basekernel);
            \draw[archarrow] (baseopt) -- (rightbus |- baseopt);
            \draw[archarrow] (baseloss) -- (rightbus |- baseloss);
            \draw[archarrow] (baseagg) -- (rightbus |- baseagg);

            \draw[archarrow] (rightbus) -- (output);

        \end{tikzpicture}
    }

    \vspace{0.18cm}

    \begin{block}{Implementation Idea}
        GradientCOBRA resolves each parameter through a factory, creates concrete modules, and keeps all modules extensible through base interfaces.
    \end{block}

\end{frame}

\begin{frame}[t]{COBRA-Based Usage}
    \vspace{-0.5cm}
    \begin{table}
        \scriptsize
        \renewcommand{\arraystretch}{1.2}
        \setlength{\tabcolsep}{1.8pt}
        \begin{tabularx}{\textwidth}{
                >{\raggedright\arraybackslash}m{2.55cm}
                >{\raggedright\arraybackslash}m{2.55cm}
                >{\raggedright\arraybackslash}X}
            \toprule
            \textbf{Parameter}          &
            \textbf{Required / Default} &
            \textbf{Purpose}                                                                          \\
            \midrule
            \texttt{estimators}         &
            \texttt{None}               &
            Base regressors; falls back to default if \texttt{None}.                                  \\
            \texttt{distance}           & \texttt{"euclidean"}      & Spatial closeness metric.       \\
            \texttt{kernel}             & \texttt{"rbf"}            & Proximity weighing engine.      \\
            \texttt{aggregator}         & \texttt{"weighted\_mean"} & Multi-expert combination logic. \\
            \texttt{loss}               & \texttt{"mse"}            & Optimization objective target.  \\
            \texttt{optimizer}          & \texttt{"grid"}           & Parameter search routine.       \\
            \texttt{learning\_rate}     & \texttt{0.1}              & Gradient adjustment step size.  \\
            \texttt{max\_iter}          & \texttt{300}              & Cap on optimizer updates.       \\
            \texttt{n\_cv}              & \texttt{5}                & Cross-validation folds.         \\
            \texttt{n\_jobs}            & \texttt{-1}               & CPU core allocation profile.    \\
            \bottomrule
        \end{tabularx}
    \end{table}
\end{frame}

\begin{frame}[fragile, t]{COBRA-Based Usage}
    \begin{minted}[
        fontsize=\scriptsize,
        breaklines,
        linenos,
        frame=single,
        framesep=2mm
    ]{python}
        from kfc_procedure.cobra import GradientCOBRA

        model = GradientCOBRA(
            estimators=["linear_regression", "random_forest_regressor"],
            kernel="rbf",           # weighting function
            distance="euclidean",   # prediction-space similarity
            loss="mse",             # optimization objective
            optimizer="grid",       # search or gradient
            max_iter=300,           # maximum iterations for optimization           
            n_cv=5,                 # cross-validation folds for hyperparameter tuning
            n_jobs=-1,              # CPU core allocation for parallel processing
            learning_rate=0.1,      # step size for gradient updates
            random_state=42,        # reproducibility random seed
        )
        model.fit(X_train, y_train)
        preds = model.predict(X_test)

    \end{minted}
\end{frame}

\subsection{Hybrid Workflow}

\begin{frame}[fragile,t]{Hybrid KFCProcedure and GradientCOBRA Workflow}
    \centering
    \scriptsize
    \vspace{-0.25cm}

    \resizebox{0.96\textwidth}{!}{%
        \begin{tikzpicture}

            \node[mainbox] (data) at (0,0) {
                Input Data\\
                $X, y$
            };

            \node[stepbox] (kstep) at (3.3,1.15) {
                K-Step\\
                Clustering
            };

            \node[stepbox] (fstep) at (6.6,1.15) {
                F-Step\\
                Local Models
            };

            \node[stepbox] (localpred) at (9.9,1.15) {
                Local\\
                Predictions
            };

            \node[stepbox] (cstep) at (13.2,1.15) {
                C-Step\\
                Combiner
            };

            \node[modulebox] (gcobra) at (6.6,-1.25) {
                GradientCOBRA\\
                Module
            };

            \node[modulebox] (predspace) at (9.9,-1.25) {
                Prediction\\
                Space
            };

            \node[modulebox] (similarity) at (13.2,-1.25) {
                Kernel / Distance\\
                Similarity
            };

            \node[modulebox] (optimize) at (16.5,-1.25) {
                Optimized\\
                Aggregation
            };

            \node[outputbox] (final) at (16.5,1.15) {
                Final Prediction\\
                \texttt{predict()}
            };

            % KFC pipeline
            \draw[archarrow] (data) -- (kstep);
            \draw[archarrow] (kstep) -- (fstep);
            \draw[archarrow] (fstep) -- (localpred);
            \draw[archarrow] (localpred) -- (cstep);
            \draw[archarrow] (cstep) -- (final);

            % COBRA plug-in
            \draw[archdashed] (cstep) --
            node[midway, right, font=\scriptsize, fill=white, inner sep=1pt] {plug-in}
            (gcobra);

            \draw[archarrow] (gcobra) -- (predspace);
            \draw[archarrow] (predspace) -- (similarity);
            \draw[archarrow] (similarity) -- (optimize);
            \draw[archarrow] (optimize) -- (final);

            \begin{scope}[on background layer]
                \node[archgroup, fit=(kstep)(fstep)(localpred)(cstep),
                    label={[grouplabel]above:KFCProcedure Pipeline}] {};

                \node[archgroup, fit=(gcobra)(predspace)(similarity)(optimize),
                    label={[grouplabel]below:COBRA-Based Aggregation Module}] {};
            \end{scope}

        \end{tikzpicture}
    }

    \vspace{0.20cm}

    \begin{block}{Hybrid Idea}
        KFCProcedure performs clusterwise learning, while GradientCOBRA can be plugged into the C-step to aggregate local predictions using prediction-space similarity.
    \end{block}

\end{frame}

\subsection{Modular and Extensible Design}

\begin{frame}[fragile,t]{Modular and Extensible Design}
    \footnotesize
    \vspace{-0.25cm}

    \begin{columns}[T,onlytextwidth]

        \column{0.43\textwidth}
        \textbf{Engineering Purpose}
        \vspace{0.12cm}

        \begin{itemize}
            \setlength\itemsep{0.18cm}

            \item \textbf{Configurable C-step:}
                  aggregation methods can be changed without changing the full KFCProcedure workflow.

            \item \textbf{Plug-in design:}
                  new combiners can be added and registered by name.

            \item \textbf{Reusable package structure:}
                  clustering, local modeling, and aggregation are separated into modules.

            \item \textbf{Maintainability:}
                  custom logic stays outside the core pipeline, reducing future modification risk.
        \end{itemize}

        \column{0.55\textwidth}
        \textbf{Example: Adding a New C-Step Combiner}
        \vspace{0.08cm}

        \begin{minted}[
            fontsize=\tiny,
            breaklines,
            frame=single,
            framesep=0.3mm,
            xleftmargin=0pt,
            xrightmargin=0pt,
            autogobble
        ]{python}
            from kfc_procedure import KFCRegressor
            from kfc_procedure.combiners import BaseCombiner, CombinerFactory
            import numpy as np

            @CombinerFactory.register("median", categories={"regression"})
            class MedianCombiner(BaseCombiner):
                def fit(self, X, y=None):
                    return self

                def combine(self, X):
                    return np.median(X, axis=1)

            model = KFCRegressor(
                divergences=["euclidean", "gkl"],
                local_model="linear_regression",
                combiner="median"
            )
        \end{minted}

    \end{columns}

    \vspace{0.12cm}

    \begin{block}{Key Benefit}
        \scriptsize
        New aggregation strategies can be added as plug-ins, while the main KFCProcedure pipeline remains unchanged.
    \end{block}

\end{frame}